% =====================================================================
% SHIFT-SURV — riscrittura §V-B-1 (I3LUNG)
% 2026-08-19
%
% Sostituisce tutto da "\subsubsection{I3LUNG}" fino a (esclusa)
% "\subsubsection{Public breast-cancer registry}".
%
% Numeri verificati contro outputs/tables/i3lung_meta.csv e
% i3lung_loco_meta.csv il 2026-08-19.
% UNICO da controllare: il totale pazienti (vedi nota [N] sotto).
% =====================================================================

\subsubsection{I3LUNG}
\label{sec:I3LUNG}

I3LUNG comprises six centers and $2074$ patients, with center sizes ranging
from $191$ to $668$.% [N] verificato: 1972 + 102 esclusi = 2074.
Each patient is described by $11$ routinely available clinical variables
selected by the treating physicians, who also chose $6$-month overall survival
as the endpoint for identifying short survivors. The fixed-horizon analysis
includes the $1972$ patients whose status at that horizon is known.

\textbf{Center-level heterogeneity.}
The centers differ both in covariate distribution and in event rate. We
quantify the first with pairwise RBF-MMD, the maximum mean discrepancy under a
Gaussian radial-basis-function kernel, where larger values indicate more
dissimilar marginal covariate distributions~\cite{gretton2012kernel}; across the
$15$ center pairs it ranges from $0.001$ to $0.160$. The $6$-month event rate
ranges from $0.16$ to $0.42$ ($\chi^2$, $p<10^{-3}$), the range the synthetic
configuration was set to reproduce (Section~\ref{sec:setup}).

These are the only two mechanisms that unlabeled data can identify, and they are
descriptive: they establish that the centers differ, not that covariate or
label-prior shift operates in isolation. The remaining four mechanisms leave no
signature computable without target labels. Neither reference model is available
here either---the latent factors are not observed, and an observed-data oracle
would require exactly the target labels that deployment withholds. The real
cohorts can therefore establish whether transport fails, but not how much of the
failure was recoverable. Supplying that decomposition is the role of the
controlled benchmark.

\textbf{Protocol.}
We evaluate the two protocols of Section~\ref{sec:setup}: adaptation over all
$30$ ordered source--target pairs, and domain generalization holding out each
center in turn with the other five pooled. Recalibration is applied after
fitting and changes only the predicted probabilities; target outcomes are used
only for evaluation. Figure~\ref{fig:realmain} reports discrimination as
$\Delta_{\mathrm{AUROC}}=\mathrm{AUROC}_{\mathrm{method}}-\mathrm{AUROC}_{\mathrm{ERM}}$
and calibration as
$\Delta_{\mathrm{ECE}}=\mathrm{ECE}_{\mathrm{ERM}}-\mathrm{ECE}_{\mathrm{method}}$,
so that positive values favor the method in both panels. For adaptation we
average the five incoming transfers within each target center, giving six
estimates whose standard errors come from the joint patient-level bootstrap;
domain generalization gives one estimate per held-out center. The six estimates
are pooled per protocol, method, and endpoint with a DerSimonian--Laird
random-effects model.

\textbf{Discrimination is not recovered.}
No method produces a practically relevant improvement on either protocol. Under
adaptation, DANN, CDAN, and IW-ERM pool \emph{below} ERM ($-0.034$, $-0.028$,
and $-0.012$ AUROC), while CORAL is indistinguishable from it ($-0.000$, $95\%$
CI $[-0.005,+0.005]$). Even the optimistic summary that selects the best method
separately for each target center has an upper $95\%$ confidence limit of
$+0.0055$, roughly a quarter of the pre-specified $0.02$ margin. Under domain
generalization, IRM and GroupDRO change AUROC by $+0.002$ and $+0.000$ with
intervals covering zero. The best-per-fold summary does reach nominal
significance here ($+0.0072$, $95\%$ CI $[+0.0006,+0.0139]$, $p=0.033$), but it
is an optimistic selection made on target performance and remains below a third
of the margin; we read it as statistically detectable and not clinically
relevant.

\textbf{Calibration fails in the direction the benchmark predicts.}
Calibration behaves differently, and the difference is mechanism-specific rather
than a matter of degree. Platt scaling improves pooled ECE by $0.018$ ($95\%$ CI
$[+0.007,+0.028]$, $p=0.001$), the only consistent gain of any method on either
endpoint. Prior adjustment moves in the opposite direction, worsening pooled ECE
by $0.074$ with extreme heterogeneity ($I^2=91\%$): it degrades calibration in
five of the six target centers, by as much as $0.16$ ECE, and improves it in
one. This is precisely the behaviour the controlled benchmark attributes to a
violated assumption. Prior adjustment recovers calibration when prevalence is
the only quantity that moves, and harms it under every other mechanism
(Fig.~\ref{fig:cal}). The I3LUNG centers differ in covariates as well as in
event rate, so the label-shift assumption does not hold and the method fails in
the anticipated direction. What the synthetic experiment identifies as an
assumption violation is observable, untreated, in a real consortium cohort.

Both patterns reproduce the synthetic result: discrimination is not recovered by
any method, while calibration changes depend on whether a method's assumption is
satisfied rather than on how much shift is present.
